\documentclass[12pt,a4paper]{article}

\usepackage[margin=2.5cm]{geometry}
\usepackage{hyperref}
\usepackage{booktabs}
\usepackage{longtable}
\usepackage{array}
\usepackage{xcolor}
\usepackage{titlesec}
\usepackage{enumitem}
\usepackage{fontenc}
\usepackage{inputenc}
\usepackage{microtype}
\usepackage{parskip}
\usepackage{fancyhdr}

% Colours
\definecolor{accent}{RGB}{30, 80, 160}
\definecolor{milestonegreen}{RGB}{34, 139, 34}
\definecolor{tablegray}{RGB}{245, 245, 245}

% Section formatting
\titleformat{\section}{\large\bfseries\color{accent}}{}{0em}{}[\titlerule]
\titleformat{\subsection}{\normalsize\bfseries}{}{0em}{}

% Header / footer
\pagestyle{fancy}
\fancyhf{}
\rhead{\textcolor{gray}{\small TestFlow — Pricing Financial Plan}}
\lhead{\textcolor{gray}{\small 2026-05-28}}
\rfoot{\textcolor{gray}{\small Page \thepage}}
\renewcommand{\headrulewidth}{0.4pt}

% Table column types
\newcolumntype{L}[1]{>{\raggedright\arraybackslash}p{#1}}
\newcolumntype{R}[1]{>{\raggedleft\arraybackslash}p{#1}}
\newcolumntype{C}[1]{>{\centering\arraybackslash}p{#1}}

\begin{document}

% ── Title block ──────────────────────────────────────────────────────────────
\begin{center}
  {\LARGE\bfseries TestFlow Pricing Financial Plan}\\[0.5em]
  {\large\color{accent} Design Specification}\\[1.5em]
  \begin{tabular}{ll}
    \textbf{Date}    & 2026-05-28 \\
    \textbf{Author}  & Parakh Sharma (solo founder) \\
    \textbf{Status}  & Approved \\
    \textbf{Related} & \texttt{PRICING.md}, \texttt{SALES\_TARGETS.md} \\
  \end{tabular}
\end{center}
\vspace{1em}
\hrule
\vspace{1.5em}

% ── Overview ─────────────────────────────────────────────────────────────────
\section{Overview}

This document models unit economics, infrastructure costs, break-even thresholds, and 36-month ARR scenarios for TestFlow (Optimus Testing). It is the financial companion to the pricing strategy in \texttt{PRICING.md}. Numbers should be reviewed every quarter and updated when actuals become available.

\textbf{Pricing model recap:} Per-workspace flat annual fee.
\begin{itemize}[noitemsep]
  \item Team --- \rupee2,40,000/year
  \item Business --- \rupee6,00,000/year
  \item Enterprise --- custom
\end{itemize}
No per-user or per-test metering. Annual upfront payment is the default.

\newcommand{\rupee}{\text{\texorpdfstring{₹}{\string₹}}}

% ── Section 1: Cost Baseline ─────────────────────────────────────────────────
\section{Cost Baseline --- Three Phases}

Costs are structured in phases keyed to customer count, not calendar time. Each phase boundary triggers a meaningful spend jump.

\subsection{Phase 0 --- Pre-revenue (current state)}

\begin{tabular}{@{}L{7cm}R{4cm}L{4.5cm}@{}}
  \toprule
  \textbf{Line item} & \textbf{Monthly cost} & \textbf{Notes} \\
  \midrule
  Supabase           & Free tier             & — \\
  Vercel             & Hobby (free)          & — \\
  EAS (Expo)         & Free                  & Limited builds \\
  Email              & Supabase built-in     & Free \\
  Claude (dev)       & ₹1,700                & Pro plan \\
  Domain (amortised) & ₹85                   & — \\
  Sentry             & Free tier             & — \\
  \midrule
  \textbf{Total}     & \textbf{{\raise.17ex\hbox{$\scriptstyle\sim$}}₹1,800/mo} & — \\
  \bottomrule
\end{tabular}

\subsection{Phase 1 --- First through \texttildelow4 customers}

Infra is upgraded on the day the first payment lands. Treat this as a one-time step-up, not a gradual ramp.

\begin{tabular}{@{}L{5.5cm}R{2.8cm}L{6cm}@{}}
  \toprule
  \textbf{Line item} & \textbf{Monthly cost} & \textbf{Notes} \\
  \midrule
  Supabase Pro          & ₹2,100  & \$25/mo · 8 GB DB, 100 K MAU, daily backups \\
  Vercel Pro            & ₹1,700  & \$20/mo · commercial domains, more build minutes \\
  EAS Production        & ₹8,300  & \$99/mo · unlimited builds, OTA updates \\
  Email                 & ₹0      & Included in Supabase Pro SMTP \\
  Claude Max 5×         & ₹11,000 & \$100/mo · upgrade from Pro \\
  Domain (amortised)    & ₹85     & — \\
  Sentry                & ₹0      & Free tier \\
  \midrule
  Subtotal              & ₹23,185 & — \\
  +10\% misc buffer     & ₹2,315  & Misc SaaS, tooling, Razorpay dashboard \\
  \midrule
  \textbf{Phase 1 total} & \textbf{{\raise.17ex\hbox{$\scriptstyle\sim$}}₹25,500/mo} & — \\
  \bottomrule
\end{tabular}

\subsection{Phase 2 --- 4+ customers}

Claude Max upgrade and monitoring tier are the main drivers. Supabase stays on Pro through \texttildelow15 customers; add a compute add-on (\texttildelow₹1,000/mo) at 10+ customers.

\begin{tabular}{@{}L{5.5cm}R{2.8cm}L{6cm}@{}}
  \toprule
  \textbf{Line item} & \textbf{Monthly cost} & \textbf{Notes} \\
  \midrule
  Supabase Pro          & ₹2,100  & +₹1,000 compute add-on at 10+ customers \\
  Vercel Pro            & ₹1,700  & — \\
  EAS Production        & ₹8,300  & — \\
  Email (Resend paid)   & ₹1,700  & \$20/mo · when invite volume exceeds Supabase SMTP \\
  Claude Max 20×        & ₹20,000 & \$200/mo · unlocked after 4--5 paying customers \\
  Domain                & ₹85     & — \\
  Sentry Team           & ₹2,200  & \$26/mo · required for production error tracking \\
  \midrule
  Subtotal              & ₹36,085 & — \\
  +10\% misc buffer     & ₹3,615  & — \\
  \midrule
  \textbf{Phase 2 total} & \textbf{{\raise.17ex\hbox{$\scriptstyle\sim$}}₹40,000/mo} & Rises to ₹43,000 with Supabase compute add-on \\
  \bottomrule
\end{tabular}

\subsection{Phase 3 --- 15+ customers (planning horizon)}

At this scale consider the Supabase Team plan (\$599/mo~$\approx$~₹50,000) for SOC\,2 compliance evidence required by enterprise buyers. First support/sales hire likely needed. Budget ₹80,000--₹1,20,000/mo total opex before payroll.

% ── Section 2: Unit Economics ─────────────────────────────────────────────────
\section{Unit Economics}

\subsection{Revenue per customer}

\begin{tabular}{@{}L{3.5cm}R{3cm}R{3cm}R{4cm}@{}}
  \toprule
  \textbf{Tier} & \textbf{Annual fee} & \textbf{MRR equiv.} & \textbf{Cash on signing} \\
  \midrule
  Team       & ₹2,40,000   & ₹20,000    & ₹2,40,000 \\
  Business   & ₹6,00,000   & ₹50,000    & ₹6,00,000 \\
  Enterprise & ₹12,00,000+ & ₹1,00,000+ & negotiated \\
  \bottomrule
\end{tabular}

\vspace{0.5em}
\textbf{Cash timing note:} Payments are annual upfront. A new Team customer means ₹2.4~L cash arrives on day~1 — not ₹20~K/month. The ``MRR equivalent'' column is for ARR tracking only; it is not a cash flow figure.

\subsection{Variable COGS per customer per month (Team tier)}

\begin{tabular}{@{}L{7cm}R{3.5cm}L{4.5cm}@{}}
  \toprule
  \textbf{Cost driver} & \textbf{Per customer/mo} & \textbf{Basis} \\
  \midrule
  Razorpay 2\% fee (amortised)        & ₹400 & 2\% × ₹2.4~L ÷ 12 \\
  Supabase incremental (DB, storage)  & ₹200 & Estimated; low at this data volume \\
  Anthropic API (AI reports)          & ₹0   & PDF export is client-side; no API call \\
  Email (incremental)                 & ₹0   & Included in flat email plan \\
  \midrule
  \textbf{Total variable COGS}        & \textbf{₹600} & Team tier \\
  \bottomrule
\end{tabular}

\vspace{0.5em}
Business and Enterprise tiers carry \texttildelow₹800--1,000/mo variable COGS (higher Razorpay fee on larger contracts, more storage).

\subsection{Gross margin by tier}

\begin{tabular}{@{}L{3.5cm}R{3cm}R{3cm}R{3cm}@{}}
  \toprule
  \textbf{Tier} & \textbf{MRR equiv.} & \textbf{Variable COGS} & \textbf{Gross margin} \\
  \midrule
  Team       & ₹20,000    & ₹600   & \textbf{97.0\%} \\
  Business   & ₹50,000    & ₹800   & \textbf{98.4\%} \\
  Enterprise & ₹1,00,000  & ₹1,000 & \textbf{99.0\%} \\
  \bottomrule
\end{tabular}

\vspace{0.5em}
These are \emph{contribution margins} — each additional customer's contribution before fixed costs. Fixed costs are covered by the break-even table in Section~\ref{sec:breakeven}.

% ── Section 3: Break-even ─────────────────────────────────────────────────────
\section{Break-even \& Milestone Table}
\label{sec:breakeven}

All figures assume Team tier (₹20~K MRR equivalent) for conservatism. Mixed Business/Enterprise deals improve the numbers.

\begin{longtable}{@{}R{2.5cm}R{2.8cm}R{2.8cm}R{2.8cm}L{4.5cm}@{}}
  \toprule
  \textbf{Customers} & \textbf{MRR equiv.} & \textbf{Phase costs} & \textbf{Monthly net} & \textbf{Milestone} \\
  \midrule
  \endfirsthead
  \toprule
  \textbf{Customers} & \textbf{MRR equiv.} & \textbf{Phase costs} & \textbf{Monthly net} & \textbf{Milestone} \\
  \midrule
  \endhead
  0  & ₹0          & ₹1,800   & $-$₹1,800             & Pre-revenue baseline \\
  1  & ₹20,000     & ₹25,500  & $-$₹5,500             & Phase~1 step-up; first month net negative \\
  2  & ₹40,000     & ₹25,500  & \textcolor{milestonegreen}{+₹14,500}  & \textcolor{milestonegreen}{✓ Default alive} \\
  3  & ₹60,000     & ₹25,500  & \textcolor{milestonegreen}{+₹34,500}  & Positive cash flow \\
  5  & ₹1,00,000   & ₹40,000  & \textcolor{milestonegreen}{+₹60,000}  & \textcolor{milestonegreen}{✓ Ramen profitable} \\
  8  & ₹1,60,000   & ₹40,000  & \textcolor{milestonegreen}{+₹1,20,000} & Solid income, no financial pressure \\
  10 & ₹2,00,000   & ₹43,000  & \textcolor{milestonegreen}{+₹1,57,000} & \textcolor{milestonegreen}{✓ GST registration threshold} \\
  15 & ₹3,00,000   & ₹43,000  & \textcolor{milestonegreen}{+₹2,57,000} & ₹36~L ARR --- Year~1 goal \\
  20 & ₹4,00,000   & ₹48,000  & \textcolor{milestonegreen}{+₹3,52,000} & Strong PMF signal; begin hiring plan \\
  30 & ₹6,00,000   & ₹55,000  & \textcolor{milestonegreen}{+₹5,45,000} & \textcolor{milestonegreen}{✓ First hire signal} \\
  50 & ₹10,00,000  & ₹65,000  & \textcolor{milestonegreen}{+₹9,35,000} & ₹1.2~Cr ARR --- pre-seed territory \\
  \bottomrule
\end{longtable}

\textbf{Key insight:} Break-even requires only 2 customers. The business becomes default-alive very quickly. The primary risk is not unit economics --- it is sales cycle length.

% ── Section 4: Scenario Model ─────────────────────────────────────────────────
\section{Three-Scenario ARR Model (36 months)}

\subsection{Scenario assumptions}

\begin{tabular}{@{}L{5cm}L{3.5cm}L{3.5cm}L{3.5cm}@{}}
  \toprule
  & \textbf{Bear} & \textbf{Base} & \textbf{Bull} \\
  \midrule
  Driver
    & Long enterprise cycles, few pilot conversions
    & \texttt{PRICING.md} Year~1 goal achieved and continued
    & Tier~2 buyer mandates TestFlow in contractor RFQs \\[0.5em]
  Pilot conversion rate & 30\% & 50\% & 70\% \\
  Avg new customers/mo  & 0.5  & 1.5  & 3.5  \\
  Annual churn          & 10\% & 5\%  & 3\%  \\
  \bottomrule
\end{tabular}

\subsection{ARR projection}

\begin{tabular}{@{}R{1.5cm}R{2cm}R{2.5cm}R{2cm}R{2.5cm}R{2cm}R{2.5cm}@{}}
  \toprule
  \textbf{Month}
    & \textbf{Bear} \newline \textbf{cust.}
    & \textbf{Bear ARR}
    & \textbf{Base} \newline \textbf{cust.}
    & \textbf{Base ARR}
    & \textbf{Bull} \newline \textbf{cust.}
    & \textbf{Bull ARR} \\
  \midrule
  3  & 0  & ₹0      & 1  & ₹2.4~L    & 3  & ₹7.2~L   \\
  6  & 1  & ₹2.4~L  & 3  & ₹7.2~L    & 6  & ₹14.4~L  \\
  12 & 3  & ₹7.2~L  & 8  & ₹19.2~L   & 20 & ₹48~L    \\
  18 & 6  & ₹14.4~L & 15 & ₹36~L     & 35 & ₹84~L    \\
  24 & 10 & ₹24~L   & 25 & ₹60~L     & 55 & ₹1.32~Cr \\
  30 & 15 & ₹36~L   & 38 & ₹91.2~L   & 78 & ₹1.87~Cr \\
  36 & 20 & ₹48~L   & 50 & ₹1.2~Cr   & 100& ₹2.5~Cr  \\
  \bottomrule
\end{tabular}

\subsection{What moves you between scenarios}

\begin{itemize}[noitemsep]
  \item \textbf{Bear $\to$ Base:} Converting one pilot every 6 weeks instead of every 3 months. Requires a sharper demo and a reference customer to cite.
  \item \textbf{Base $\to$ Bull:} A single Tier~2 enterprise (e.g.\ Adani, Tata Power) mandates TestFlow in their contractor RFQs. One such win can add 10--20 workspaces at once. This is the asymmetric upside.
\end{itemize}

% ── Section 5: GST & Tax ─────────────────────────────────────────────────────
\section{GST \& Tax Notes}

\begin{itemize}[noitemsep]
  \item \textbf{Registration threshold:} ₹20~L cumulative revenue in a financial year. At Team tier, this is \texttildelow10 customers (10~×~₹2.4~L~=~₹24~L). Register proactively 30 days before crossing this.
  \item \textbf{GST rate:} 18\% on software/SaaS services. Invoice as ₹2,40,000 + ₹43,200 GST = ₹2,83,200 total. The ₹43,200 is collected and remitted --- it is not your revenue.
  \item \textbf{Input tax credit (ITC):} Supabase, Vercel, and EAS are billed in USD by US entities and likely carry no ITC. Claude Max billed via Anthropic US --- no ITC. Budget these costs as gross, not net of ITC.
  \item \textbf{Razorpay GST:} Razorpay's 2\% fee attracts 18\% GST on the fee itself (effective rate: 2.36\%). Adjust cost estimates upward by \texttildelow0.36\% of each transaction value.
\end{itemize}

% ── Section 6: Key Metrics ────────────────────────────────────────────────────
\section{Key Metrics to Track Monthly}

\begin{tabular}{@{}L{4.5cm}L{3.5cm}L{7cm}@{}}
  \toprule
  \textbf{Metric} & \textbf{Target} & \textbf{What it tells you} \\
  \midrule
  ARR per workspace
    & $\geq$~₹2~L
    & Discount discipline. Below ₹1.4~L = stop discounting (per \texttt{PRICING.md}). \\[0.4em]
  Monthly burn
    & vs.\ phase threshold
    & Phase~1: ₹25.5~K · Phase~2: ₹40~K. Creep above threshold without customer growth = investigate. \\[0.4em]
  Active pilots
    & $\leq$ 5 at once
    & Per \texttt{PRICING.md} cap. More than 5 concurrent = distraction. \\[0.4em]
  Pilot → paid conversion
    & $\geq$ 50\%
    & Below 40\% = product/onboarding problem, not a pricing problem. \\[0.4em]
  Net Revenue Retention
    & $\geq$ 100\%
    & Track from Year~2. NRR~>~100\% means expansions offset churn. \\[0.4em]
  Cash in bank
    & $\geq$ 6 months Phase~2 costs
    & Minimum buffer: 6~×~₹40~K~=~₹2.4~L. One annual signing can refill this immediately. \\
  \bottomrule
\end{tabular}

% ── Section 7: Financial Milestones ──────────────────────────────────────────
\section{Financial Milestones Summary}

\begin{longtable}{@{}L{4.5cm}L{4cm}L{6.5cm}@{}}
  \toprule
  \textbf{Milestone} & \textbf{Trigger} & \textbf{Action} \\
  \midrule
  \endfirsthead
  \toprule
  \textbf{Milestone} & \textbf{Trigger} & \textbf{Action} \\
  \midrule
  \endhead
  Upgrade infra        & 1st payment received   & Supabase Pro + Vercel Pro + EAS Production on the same day \\[0.3em]
  Default alive        & 2nd customer pays      & Business survives indefinitely on current trajectory \\[0.3em]
  Ramen profitable     & 5th customer pays      & Founder draws ₹60~K/mo personal income \\[0.3em]
  GST registration     & \texttildelow10th customer or ₹20~L ARR & File 30 days before crossing threshold \\[0.3em]
  Claude Max 20×       & 4th--5th customer      & Reinvest margin into faster product velocity \\[0.3em]
  Sentry Team          & 4th--5th customer      & Production error monitoring for enterprise conversations \\[0.3em]
  First marketing spend & 8th customer          & Budget ₹20~K/mo for LinkedIn Sales Navigator + outreach tooling \\[0.3em]
  First hire evaluation & 20--25 customers      & Sales + onboarding support; budget ₹80--₹1.2~L/mo \\[0.3em]
  Pricing review        & 12 months from first paid customer & Raise Team to ₹3~L/yr per \texttt{PRICING.md} Year~2 plan \\
  \bottomrule
\end{longtable}

% ── Appendix ─────────────────────────────────────────────────────────────────
\appendix
\section{Infra Upgrade Costs (one-time on first customer)}

\begin{tabular}{@{}L{3.5cm}L{6cm}R{4cm}@{}}
  \toprule
  \textbf{Item} & \textbf{One-time action} & \textbf{Recurring from that month} \\
  \midrule
  Supabase      & Upgrade plan in dashboard         & ₹2,100/mo \\
  Vercel        & Upgrade plan in dashboard         & ₹1,700/mo \\
  EAS           & Enable Production plan            & ₹8,300/mo \\
  Claude Max 5× & Change subscription               & ₹11,000/mo \\
  Email         & Configure Supabase SMTP or Resend & ₹0--₹1,700/mo \\
  \midrule
  \textbf{Total step-up} & — & \textbf{+₹23,700/mo vs.\ Phase~0} \\
  \bottomrule
\end{tabular}

\vspace{0.5em}
This step-up is covered by customer~2's MRR equivalent within the same month.

\end{document}
